\section{Related literature and the scope of the contribution}
\label{sec:extended-literature}

This supplement distinguishes the paper's empirical question from established results in predictive execution, uncertainty-aware trading and decision-focused learning. The study follows the author's execution preprint \citep{lai2026}, but its contribution must be assessed through the additional forecast comparison, chronological experiments and reproducible diagnostics. Neither the cryptocurrency setting nor the use of signals in execution is itself new. The relevant question is whether a forecast that predicts a terminal return also identifies price differences across the times at which a fixed order can actually be executed.

\subsection{Predictive execution and the temporal shape of forecasts}

The impact--risk trade-off in liquidation is foundational \citep{almgren2001}. Subsequent work incorporates information about order flow and future prices into trading controls. \citet{cartea2016} show how weighted expected future net order flow changes an execution strategy. \citet{bechler2015} incorporate dynamic order-flow imbalance and consider an endogenous execution horizon, while \citet{cartea2018} investigate trading decisions informed by order-book signals. These approaches already make the timing of information economically relevant. They cannot be represented adequately by a general claim that earlier execution research ignores predictive information.

Several papers are especially close to the temporal-profile argument. The deterministic, time-zero-conditioned formulation in \citet{lehalle2019} uses the integrated expected signal across execution times and removes a fixed initial-price term from the objective. \citet{forde2022} formulate execution with Gaussian signals and power-law resilience using a conditional remaining price change. The general-propagator formulation of \citet{abijaber2025} likewise contains conditional price contrasts relative to the terminal price. Thus, the fact that relative expected prices matter for allocating a fixed quantity across time is already embedded in established execution objectives. The centered-profile identity used here is an elementary finite-grid representation of that fact, rather than a new principle of optimal execution.

The distinction is also related to the persistence of predictors. In dynamic portfolio choice with transaction costs, \citet{garleanu2013} derive adjustment toward an aim portfolio that depends on future expected targets. That setting differs from compulsory fixed-volume liquidation, but it establishes why a predictor's temporal behavior can affect its trading value. Learning latent alpha states \citep{casgrain2019}, adapting execution to signals under temporary and transient impact \citep{neuman2022}, and responding to imminent trade signals under liquidity-dependent impact \citep{bank2026} provide further examples of richer information structures than a single terminal-direction forecast.

Our narrower diagnostic holds the feasible set and other cost terms fixed. A common additive shift in all execution-slot price forecasts contributes the same amount to every equal-volume schedule. It therefore cannot alter their linear price comparison. A constant drift is different: it generally changes expected prices unevenly across execution times. Nor does level cancellation establish invariance when forecasts change constraints, impact parameters or the normalization of a reported performance ratio. The empirical comparison concerns the centered temporal profile under the stated replay assumptions; it does not establish a general dominance result over these predictive-control models.

\subsection{Uncertainty, shrinkage and fallback policies}

Uncertainty-aware trading is also an established direction. \citet{cartea2017} study model uncertainty in algorithmic trading, primarily in a market-making setting. Liquidation under ambiguity aversion is addressed directly in \citet{cartea_liquidation2017}, and \citet{horst2022} study liquidation under factor uncertainty. \citet{demarch2018} examine uncertainty associated with threshold-based signal controls. These are distinct formulations, rather than interchangeable guarantees for a fitted scalar attenuation rule. Portfolio parameter uncertainty \citep{kan2007} supplies a related motivation for cautious estimation, but its portfolio-choice conclusions do not automatically transfer to an execution schedule.

The closest recent uncertainty-gating comparison is \citet{irshad2026}. Their execution framework gates next-interval forecast-driven action using uncertainty estimates and otherwise uses a VWAP fallback. Consequently, ignoring an uncertain forecast or retaining a baseline is not the novelty claimed here. Their equity-market experiment also cannot validate a cryptocurrency replay. The relevant differences are the quantity being diagnosed, the temporal forecast representation, the calibration procedure and the evidence produced under our own implementation.

Errors in execution-cost models require separate attention. \citet{hey2023} analyze costs of misspecified price impact, while \citet{neumanoffline2023} propose pessimistic offline learning for propagator models. Stochastic impact enters the control problem in \citet{fouque2022}; \citet{macri2024} investigate reinforcement learning when liquidity varies over time. These studies motivate treating a fixed impact coefficient as an assumption rather than a measured market property. Baseline-constrained policy improvement is also established in reinforcement learning \citep{laroche2019}, but its safety results do not constitute guarantees for our calibration or resampling procedure. Here, attenuation and fallback are evaluated empirically within the specified replay.

\subsection{Prediction quality, decision quality and coverage}

The difference between statistical prediction accuracy and downstream decision quality is central to smart predict-then-optimize learning \citep{elmachtoub2022}. \citet{liu2021} establish surrogate-risk and calibration results under particular assumptions on the optimization problem. The survey and benchmark of \citet{mandi2024} place such methods within the broader decision-focused learning literature, and \citet{gupta2024} develop a directional-gradient approach. These works precede the motivation for evaluating forecasts through their consequences for a constrained decision. The present diagnostic is task-specific and does not introduce a general decision-focused learning framework or inherit its theoretical guarantees.

Temporal dependence adds a further distinction. \citet{liudependent2024} investigate predict-then-optimize risk bounds with dependent data, illustrating why independence-based conclusions require additional work in time series. Prediction-interval coverage is likewise different from an execution-cost guarantee. Adaptive conformal inference targets coverage under distribution shift \citep{gibbs2021}; departures from exchangeability are explicitly analyzed by \citet{barber2023}. Time-series adaptations include \citet{zaffran2022} and \citet{xu2021}, with different constructions and conditions. These references provide context for interpreting uncertainty-aware execution methods. They do not justify calling empirical shrinkage, a bootstrap interval or a baseline fallback distribution-free protection against trading losses.

\subsection{Cryptocurrency evidence and backtest interpretation}

Cryptocurrency execution has direct precedents. \citet{bundi2024} study optimal trade execution in cryptocurrency markets, and \citet{schnaubelt2022} investigate cryptocurrency limit-order placement using reinforcement learning and detailed market data. The latter problem includes information about orders and trades that a bar-price replay cannot reproduce. \citet{barucci2023} examine market impact and efficiency in cryptoassets, while \citet{makarov2020} document exchange segmentation and arbitrage-related price differences. The microstructure review of \citet{almeida2024} supplies broader context. Together these studies motivate asset-, exchange-, horizon- and data-resolution-specific interpretation; they do not establish profitable execution from terminal-return predictability alone.

Finally, dependence and model selection are separate sources of uncertainty. The stationary bootstrap of \citet{politis1994} is a specific dependence-aware resampling construction, not a blanket justification for every block-bootstrap implementation. \citet{bailey2017} examine overfitting created by selecting among backtests, and \citet{arnott2019} advocate a disciplined empirical research protocol. Accordingly, chronological splits, calibration choices, inspected periods and subsequent design changes should be disclosed. An expanded study designed after an initial result has been observed remains exploratory in that respect. The contribution is a transparent comparison of forecast targets and execution consequences under stated assumptions, with uncertainty and negative findings reported alongside any apparent gains.
